% Gesamttest «Wärme» — Quelle des PDFs (python3 scripts/build-lp-pdf.py waerme).
% Musterlösung und Raster: bewertungspaket.tex. Alle Werte mit python3 nachgerechnet (08.10.2026).
% Fassung 1.1 nach /lp-pruefung: G2 Bilanz mit Zustandsänderung (Eis in Limonade, nicht alles schmilzt),
% G3 Heizkurve lesen und weiterzeichnen (K3), G4 Leistung rückwärts statt Nachbau der Übung «heizwert»,
% G6 Solarkocher statt Heizkörper-Problem des Clips, G7 rückwärts und ehrlich zum Einschichtmodell.
% Fassung 1.2 nach der zweiten Prüfung: G1 b Mischtemperatur (Kapitel 2, Bilanz ohne Zustandsänderung),
% G2 mit 0.25 kg Limonade (geschmolzen 0.063 kg, übrig 0.037 kg — nicht mehr zufällig gleich), G3 mit 0.40 kg
% (nicht wie Leiste sim3 A1), G4 Campingkocher mit Wirkungsgrad aus der Messung (nicht wie Übung «heizwert»),
% G7 c im Einschichtmodell mit der einen Einstellung (390, 240, 150 W/m²). Fassung 1.3 (dritte Prüfung): G4 mit 2.0 l
% von 12 °C und 30 g Butan (nicht wie das Clip-Problem «Topf»), G6 b Holzgriff statt «fühlt sich kälter an» (wie Aufgabe 6a).
% Keine Aufgabe aus uebungstest-waermelehre; geprüft auch mit scripts/lp/waerme/pruef_fest.py.
\documentclass[11pt]{article}
\usepackage{../lp-druck}
\renewcommand{\lpname}{Wärme}
\renewcommand{\lpteilgebiet}{5.2}
\renewcommand{\lpfuss}{Gesamttest · 25 Punkte}
\begin{document}
\lpkopf{Gesamttest}{%
  \vspace{4pt}\noindent\begin{tabularx}{\linewidth}{@{}X X@{}}
  Code (kein Name): \hrulefill & Punkte: \hrulefill\ / 25
  \end{tabularx}}

\begin{lpinfo}{Anleitung}
\textbf{Zeit:} rund 40 Minuten. \textbf{Hilfsmittel:} Taschenrechner und Formelsammlung, in allen Teilen.
Schreib zu jeder Aufgabe zuerst die Formel, dann die Werte \textbf{mit Einheiten}. Punkte gibt es für den Weg,
für das Ergebnis mit Einheit, fürs Ablesen und Zeichnen und für Begründungen in eigenen Worten. Rechne mit
\mbox{$c_\text{Wasser} = 4182\;\text{J/(kg·K)}$}, \mbox{$L_\text{v} = 2256\;\text{kJ/kg}$} (Wasser),
\mbox{$1\;\text{kWh} = 3.6\;\text{MJ}$} und \mbox{$\sigma = 5.67 \cdot 10^{-8}\;\text{W/(m}^2\text{K}^4)$}.
Danach mit dem \emph{Bewertungspaket} bewerten (Scan ohne Namen und Standort).
\end{lpinfo}

\lpteil{Teil A · Wärme, Bilanz und Phasenwechsel}{Kapitel 1, 2 und 3 · 12 Punkte}

\begin{lpaufgabe}{G1}{Ethanol im Wasserbad}{3 P}
Eine dünnwandige Flasche mit \(0.35\;\text{kg}\) Ethanol (\(c = 2430\;\text{J/(kg·K)}\)) von \(60\;^\circ\text{C}\) steht in \(0.50\;\text{kg}\) Wasser von \(18\;^\circ\text{C}\). Flasche, Gefäss und Umgebung nehmen nichts auf.
\begin{enumerate}[label=\alph*),leftmargin=*,itemsep=2pt,topsep=2pt]
\item Jemand sagt: «Das heisse Ethanol gibt dem Wasser so lange Temperatur ab, bis beide gleich warm sind.» Was ist an dem Satz falsch? Erkläre den Unterschied zwischen Temperatur und Wärme.
\item Welche Temperatur haben Ethanol und Wasser am Ende?
\end{enumerate}
\lpplatz{6}
\end{lpaufgabe}

\begin{lpaufgabe}{G2}{Eis in der Limonade}{4 P}
In \(0.25\;\text{kg}\) Limonade von \(20\;^\circ\text{C}\) (Stoffwerte wie Wasser) kommen \(0.10\;\text{kg}\) Eis von \(0\;^\circ\text{C}\) (\mbox{\(L_\text{f} = 334\;\text{kJ/kg}\)}). Rechne ohne Verluste.
\begin{enumerate}[label=\alph*),leftmargin=*,itemsep=2pt,topsep=2pt]
\item Wie viel Wärme kann die Limonade höchstens abgeben, bis sie \(0\;^\circ\text{C}\) hat?
\item Reicht das, um alles Eis zu schmelzen? Begründe mit einer Rechnung.
\item Welche Temperatur hat das Glas am Ende, und wie viel Eis schwimmt dann noch darin?
\end{enumerate}
\lpplatz{5}
\end{lpaufgabe}

\begin{lpaufgabe}{G3}{Eis auf der Kochplatte}{5 P}
\(0.40\;\text{kg}\) Eis von \(0\;^\circ\text{C}\) liegen in einem Topf auf einer Kochplatte, die dem Inhalt gleichmässig Energie zuführt. Das Diagramm zeigt die gemessene Temperatur; nach \(12\;\text{min}\) wurde die Messung abgebrochen.\par\smallskip
\begin{center}\begin{tikzpicture}[x=0.38cm,y=0.042cm,>=latex]
\foreach \x in {0,1,...,40} \draw[lplinie!70,very thin] (\x,-10) -- (\x,110);
\foreach \y in {-10,0,...,110} \draw[lplinie!70,very thin] (0,\y) -- (40,\y);
\foreach \y in {0,20,...,100} {\draw[lplinie] (0,\y) -- (40,\y); \node[left] at (0,\y) {\small \y};}
\foreach \x in {5,10,...,40} {\draw[lplinie] (\x,-10) -- (\x,110); \node[below] at (\x,-10) {\small \x};}
\draw[->] (0,0) -- (41.5,0);
\draw (0,-10) -- (0,0);
\draw[->] (0,-10) -- (0,118) node[right] {\small \(\vartheta\) in \(^\circ\)C};
\node[above left] at (41.5,0) {\small \(t\) in min};
\draw[very thick,lpbernstein] (0,0) -- (4,0) -- (9,100) -- (12,100);
\fill[lpbernstein] (12,100) circle (1.6pt);
\node[lpbernstein,anchor=south west,font=\small] at (12.3,101) {Messung abgebrochen};
\end{tikzpicture}\end{center}
\begin{enumerate}[label=\alph*),leftmargin=*,itemsep=2pt,topsep=2pt]
\item Lies ab, wie lange das Schmelzen dauert und wie schnell danach das Wasser wärmer wird. Berechne daraus, wie viel Leistung die Platte dem Inhalt zuführt.
\item Prüfe mit dem Diagramm und deinem Ergebnis aus a) die spezifische Schmelzwärme von Eis nach.
\item Berechne, wann alles Wasser verdampft ist, und zeichne die Kurve im Diagramm bis zu diesem Zeitpunkt kräftig weiter.
\end{enumerate}
\lpplatz{11}
\end{lpaufgabe}

\lpteil{Teil B · Brennstoffe und Energiesysteme}{Kapitel 4 und 5 · 7 Punkte}

\begin{lpaufgabe}{G4}{Der Campingkocher}{4 P}
Ein Campingkocher bringt \(2.0\;\text{l}\) Wasser in \(8.0\;\text{min}\) von \(12\;^\circ\text{C}\) zum Sieden (\(100\;^\circ\text{C}\)). Dabei verbrennt er \(30\;\text{g}\) Butan (Heizwert \(H = 45.7\;\text{MJ/kg}\)).
\begin{enumerate}[label=\alph*),leftmargin=*,itemsep=2pt,topsep=2pt]
\item Wie viel Wärme nimmt das Wasser auf?
\item Wie gross ist der Wirkungsgrad des Kochers?
\item Welche Leistung führt die Flamme zu, in Kilowatt?
\item Wohin geht der Rest der Energie? Begründe, warum er nicht verschwunden ist.
\end{enumerate}
\lpplatz{10}
\end{lpaufgabe}

\begin{lpaufgabe}{G5}{Biogas auf dem Bauernhof}{3 P}
Ein Bauernhof gewinnt aus Gülle und Mist im Jahr \(60\,000\;\text{m}^3\) Biogas mit je \(6.0\;\text{kWh}\) Energie. Ein Blockheizkraftwerk (Wärme-Kraft-Kopplung) macht daraus \(35\;\%\) Strom und \(55\;\%\) Nutzwärme.
\begin{enumerate}[label=\alph*),leftmargin=*,itemsep=2pt,topsep=2pt]
\item Wie viel Strom und wie viel Nutzwärme liefert die Anlage im Jahr?
\item Ist die Wärme-Kraft-Kopplung eine erneuerbare Energiequelle? Begründe.
\item Nenne einen Vorteil der Biogasanlage gegenüber einer Photovoltaikanlage, die im Jahr gleich viel Strom liefert. Begründe.
\end{enumerate}
\lpplatz{8}
\end{lpaufgabe}

\lpteil{Teil C · Wärmetransport und Treibhauseffekt}{Kapitel 6 und 7 · 6 Punkte}

\begin{lpaufgabe}{G6}{Der Solarkocher}{3 P}
Ein einfacher Solarkocher besteht aus einem spiegelnden Trichter, einem schwarzen Topf und einem durchsichtigen Plastiksack, der um den Topf gebunden ist. Der Topf steht auf einem Holzbrett.
\begin{enumerate}[label=\alph*),leftmargin=*,itemsep=2pt,topsep=2pt]
\item Erkläre für jeden der drei Teile (Trichter mit schwarzem Topf, Plastiksack, Holzbrett), welchen Weg der Wärme er nutzt oder bremst. Begründe.
\item Eine Bratpfanne hat einen Griff aus Holz statt aus Stahl. Begründe mit den Wärmeleitfähigkeiten (Stahl rund \(50\;\text{W/(m·K)}\), Holz rund \(0.15\;\text{W/(m·K)}\)), warum man den Griff auch bei heisser Pfanne anfassen kann.
\end{enumerate}
\lpplatz{9}
\end{lpaufgabe}

\begin{lpaufgabe}{G7}{Strahlung am Boden}{3 P}
\begin{enumerate}[label=\alph*),leftmargin=*,itemsep=2pt,topsep=2pt]
\item Erkläre, warum die Atmosphäre das Sonnenlicht anders behandelt als die Strahlung, die der Boden abgibt, und was die Treibhausgase mit der aufgenommenen Strahlung tun.
\item Ein Boden strahlt \(400\;\text{W/m}^2\) Wärmestrahlung ab (idealer Strahler). Welche Temperatur hat er, in Grad Celsius?
\item Im Einschichtmodell des Leitprogramms strahlt der Boden \(390\;\text{W/m}^2\) ab und bekommt \(240\;\text{W/m}^2\) Sonnenlicht und \(150\;\text{W/m}^2\) Gegenstrahlung. Gemessen ist die Gegenstrahlung aber rund \(340\;\text{W/m}^2\). Begründe, warum der Boden trotzdem nicht wärmer wird.
\end{enumerate}
\lpplatz{8}
\end{lpaufgabe}

\vfill
{\small\color{lpgrau}Bewerten: mit dem Bewertungspaket (Musterlösung, Punkteraster, Auftrag an eine KI) — erst nach dem Lösen öffnen.}
\end{document}
